% from ega4-20-fr.tex
\oldpage[IV-4]{231}
des éléments réguliers de $B$, et comme on a vu que $S\subset T$, l'homomorphisme $S^{-1}B\to T^{-1}B$ est injectif ; autrement dit, cela prouve que l'homomorphisme \sref{IV.20.1.11.1} $f^*(\sh{M}_X)\to\sh{M}_{X'}$ est \emph{injectif} (d'où la dernière assertion de l'énoncé).
Le quotient $f^*(\sh{M}_X)/\sh{O}_{X'}$ s'identifie à un sous-$\sh{O}_{X'}$-module de $\sh{M}_{X'}/\sh{O}_{X'}$, et $(f^*(\sh{M}_X)/\sh{O}_{X'})_{x'}$ s'identifie à $(\sh{M}_X/\sh{O}_X)_x\otimes_{\sh{O}_{X,x}}\sh{O}_{X',x'}$.
Supposons alors que $x\notin\operatorname{dom}(\vphi)$ ; l'image de $\vphi_x$ dans $(\sh{M}_X/\sh{O}_X)_x$ est donc $\neq0$ ; par fidèle platitude, on en déduit qu'il en est de même de l'image de $(\vphi\circ f)_{x'}$ dans $(\sh{M}_{X'}/\sh{O}_{X'})_{x'}$, donc $x'\notin\operatorname{dom}(\vphi\circ f)$, ce qui achève la démonstration.
\end{proof}

\begin{remark}[20.1.13]
\label{IV.20.1.13}
Soit $X$ un espace analytique complexe \emph{réduit} ; alors la notion de fonction méromorphe sur $X$ définie ci-dessus coïncide avec la notion usuelle.
Considérons d'autre part un préschéma $Y$, localement de type fini sur le corps $\bb{C}$ ; on sait alors qu'on peut associer à $Y$ un espace analytique $Y^\mathrm{an}$ ayant même espace topologique sous-jacent, et que le morphisme canonique $f:Y^\mathrm{an}\to Y$ est \emph{plat}~\cite{IV-37} ; en vertu de \sref{IV.20.1.12}, l'homomorphisme canonique $u\mapsto u\circ f$ de $M(Y)$ dans $M(Y^\mathrm{an})$ est donc partout défini et est injectif ; mais il n'est pas \emph{surjectif} en général.
Par exemple, lorsque $Y=\bb{V}_{\bb{C}}^r$ (\textbf{Err}\textsubscript{\textbf{III}},~14) est l'espace affine de dimension $r$ sur $\bb{C}$, $M(Y)$ s'identifie canoniquement au corps $R(Y)$ des fonctions rationnelles sur $Y$ \sref{IV.20.2.13}[(i)], tandis que $M(Y^\mathrm{an})$ est le corps des fonctions méromorphes usuelles sur $\bb{C}^r$.
En raison de ce fait, il est souvent préférable, en Géométrie algébrique, de s'abstenir de la terminologie introduite dans cette section, et d'utiliser la terminologie équivalente de « pseudo-fonction » qui va être définie ci-dessous.
\end{remark}

\subsection{Pseudo-morphismes et pseudo-fonctions}
\EGAReaderTarget{ega4.c04.subsec.20-2}{IV.20.2}

Les seuls espaces annelés dont il sera question dans cette section sont les \emph{préschémas}.

\begin{env}[20.2.1]
\EGAReaderTarget{ega4.c04.statement.20-2-1}{IV.20.2.1}
Rappelons \sref{IV.11.10.2} que dans un préschéma $X$, on dit qu'un ouvert $U$ est \emph{schématiquement dense} si, pour tout ouvert $V$ de $X$, l'homomorphisme canonique $\Gamma(V,\sh{O}_X)\to\Gamma(V\cap U,\sh{O}_X)$ est \emph{injectif}.

Considérons deux préschémas $X$, $Y$, deux ouverts \emph{schématiquement denses} $U$, $U'$ de $X$ ; on dit que deux morphismes $u:U\to Y$, $u':U'\to Y$ sont \emph{équivalents} s'il existe un ouvert $U''\subset U\cap U'$, \emph{schématiquement dense} dans $X$, tel que $u|U''=u'|U''$.
Comme il résulte aussitôt de la définition des ouverts schématiquement denses que l'intersection de deux tels ouverts en est un autre, il est immédiat que la relation précédente est bien une relation d'équivalence.
Une classe d'équivalence suivant cette relation s'appelle un \emph{pseudo-morphisme} de $X$ dans $Y$, ou une \emph{application rationnelle stricte} de $X$ dans $Y$.

Si $S$ est un préschéma et $X$, $Y$ deux $S$-préschémas, on appelle \emph{pseudo-$S$-morphisme} la classe d'équivalence (pour la relation précédente) d'un $S$-morphisme d'un ouvert schématiquement dense de $X$ dans $Y$.
Un pseudo-morphisme n'est donc autre chose qu'un pseudo-$S$-morphisme pour $S=\Spec(\bb{Z})$.

On note $\operatorname{Ps.hom}(X,Y)$ (resp. $\operatorname{Ps.hom}_S(X,Y)$) l'ensemble des pseudo-morphismes (resp. des pseudo-$S$-morphismes) de $X$ dans $Y$.
\end{env}

\begin{env}[20.2.2]
\EGAReaderTarget{ega4.c04.statement.20-2-2}{IV.20.2.2}
